\documentclass{beamer} [10.12.2024]
\usepackage[utf8]{inputenc}
\usepackage[ngerman]{babel}
\usepackage[T1]{fontenc}
\usepackage{lmodern}
\usepackage{array}
\usepackage{graphicx}
\usepackage{mathtools}
\usepackage{amsmath} 
\usepackage{subfig}
\usepackage{hyperref}
\usepackage{algorithm}
\usepackage{algpseudocode}

\captionsetup[figure]{labelformat=simple, labelsep=colon}
\renewcommand\thefigure{\arabic{figure}}
\setbeamertemplate{caption}[numbered]
\bibliographystyle{plain} 

\definecolor{lightblue}{rgb}{0.47, 0.47, 0.85}


\mode<presentation>{%
	\usetheme{Berlin}
	\setbeamercovered{transparent}
}


\title{CZGame}
\subtitle{Ein interaktives Point-and-Click Spiel}
\author{Jeje C., Gustav G., Rita G., Jonas G., Jonas L., Sophie P., Emil R., Antonia V., Arthur W., Ashley W.}
\institute[CZG]{Carl-Zeiss-Gymnasium Jena}
\date{\today}

\begin{document}
	
	\section{Minigames}
	\begin{frame}
	\begin{center}
		\Huge Minigames
	\end{center}
\end{frame}

\begin{frame}
	\frametitle{Euer Titel}
\end{frame}

\begin{frame}
	\frametitle{Euer Titel}
\end{frame}

\begin{frame}
	\frametitle{Euer Titel}
\end{frame}

\begin{frame}
	\frametitle{Euer Titel}
\end{frame}

\begin{frame}
	\frametitle{Euer Titel}
\end{frame}

\begin{frame}
	\frametitle{Euer Titel}
\end{frame}

\begin{frame}
	\frametitle{Euer Titel}
\end{frame}

\end{document}